\documentclass[10pt,a4paper]{amsart}
\usepackage[margin=19mm]{geometry}
\usepackage{amsmath,amssymb,mathtools}
\usepackage[scr=rsfs,cal=euler]{mathalfa}
\usepackage{xfrac}
\usepackage[unicode,hidelinks,bookmarks=false]{hyperref}
\newcommand{\ie}{i.e.\ }
\newcommand{\cf}{cf.\ }
\newcommand{\resp}{respectively\ }
\newcommand{\Subsection}{\subsection}
\providecommand{\nobreakdash}{\nobreak\hbox{-}}
\setlength{\emergencystretch}{2em}
\multlinegap0pt
\usepackage{amssymb}
\usepackage[shortlabels]{enumitem}
\usepackage{tikz}
\usepackage{mathtools}
\usepackage{xfrac}
\newtheorem{theorem}{Theorem}
\newcommand{\A}{\mathbb{A}}
\newcommand{\C}{\mathbb{C}}
\newcommand{\CP}{\mathscr{P}}
\DeclareMathOperator{\Spec}{Spec}
\title{The projective coinvariant algebra, Young invariants and \\ bigraded coordinate rings of Segre embeddings}
\author{Bal\'azs Szendr\H{o}i}
\date{Source: arXiv:2602.15017v1 (CC BY 4.0)}
\begin{document}
\maketitle

\noindent\emph{Source and licence.} The projective coinvariant algebra, Young invariants and \\ bigraded coordinate rings of Segre embeddings. Version arXiv:2602.15017v1 (CC BY 4.0), distributed by arXiv under the Creative Commons Attribution 4.0 International licence, http://creativecommons.org/licenses/by/4.0/, as recorded in the arXiv licence metadata for this exact version and shown on the abstract page https://arxiv.org/abs/2602.15017v1. Full licence notice: \texttt{licenses/arxiv-2602.15017-CC-BY-4.0.txt}.\par\smallskip

\noindent\textit{Editorial scope.} Counted unit: the article's theorem thm\_main\_def (source0.tex lines 360--366). The article's complete original proof of it is delivered with it (source0.tex lines 367--392, one \texttt{proof} environment of 26 lines). No mathematics is written, repaired or re-derived: the statement and the proof are the article's own lines, in the article's own notation, and no retained line is substituted or re-flowed.  No further source-local context is needed: every cross-reference of every retained line resolves inside the two delivered spans.  Nothing was omitted from any retained span, so the package carries no omission record.  No retained line contains a citation, so the wrapper carries no bibliography and no bibliography entry.  No retained line contains a figure environment, an image-inclusion command or drawing code, so no image is delivered and the package declares no asset dependency.  Editorial formatting only, in addition: an AMS \texttt{amsart} preamble that restates the article's own declarations and macros used by the retained lines, namely the environments \texttt{theorem} on separate counters, together with the standard packages those lines need.  The retained environments print in this document as Theorem 1 in order of appearance, whereas the article prints the same items with the numbering of its own full document.  The complete native source of the article as distributed in the arXiv e-print for this exact version is bundled unchanged as \texttt{sources/194766/source0.tex}.
\begin{theorem}\label{thm_main_def}\begin{enumerate}
\item[(i)] The inclusion $\C[s_0, s_1]\rightarrow \CP_n$ is a flat map of algebras, equivariant with respect to the action of the torus $(\C^*)^2$, as well as the symmetric group $S_n$. 
\item[(ii)] The central fibre $(\CP_n)_{(0,0)}$ is bigraded $S_n$-isomorphic to the projective coinvariant algebra $P_n$. 
\item[(iii)] Over points $(s_0, 0)$ and $(0,s_1)$ with $s_i\neq 0$, the fibres $(\CP_n)_{(s_0,0)}$, $(\CP_n)_{(0,s_1)}$ are graded $S_n$-isomorphic to the classical coinvariant algebra $R_n$.
\item[(iv)] Over points $(s_0, s_1)$ with $s_0s_1\neq 0$, the fibre $(\CP_n)_{(s_0,s_1)}$ is an (ungraded) finite-dimensional algebra of dimension $n!$, $S_n$-isomorphic to the group algebra $\C S_n$. 
\end{enumerate}
\end{theorem}

\begin{proof}

The central fibre~$(\CP_n)_{(0,0)}$ is by definition isomorphic to the projective coinvariant algebra $P_n$. On the other hand, being a regular sequence is an open property, so the corresponding elements give regular sequences in fibres near the origin, and hence, using the torus action, over the whole affine plane. 

Let us consider the fibre $(\CP_n)_{(s_0,0)}$ over a point $(s_0, 0)$ with $s_0\neq 0$. Re-scaling using the first component of the torus action, we can assume that $s_0=1$. In this algebra $(\CP_n)_{(1,0)}$, the image of the element $x_{\scriptscriptstyle \emptyset}$ gets identified with the unit. Define a $\C$-algebra homomorphism 
\[\psi_n\colon \C[u_1, \ldots, u_n]\to (\CP_n)_{(1,0)}\]
by $\psi_n(u_i)=(x_{\scriptscriptstyle\{i\}})_{(1,0)}$. Then for $i\not\in I$, using the relations $x_{\scriptscriptstyle\emptyset} x_{\scriptscriptstyle I\cup \{i\}} = x_{\scriptscriptstyle\{i\}} x_{\scriptscriptstyle I}$ inductively, we deduce 
\[\psi_n\left(\prod_{i\in I} u_i\right) = (x_{\scriptscriptstyle I})_{(1,0)}  \in (\CP_n)_{(1,0)}.\]
In particular, if $e_k(u_1,\ldots, u_n)$ denotes the elementary symmetric functions of the variables $u_1, \ldots, u_n$, we get
\[\psi_n\left(e_k(u_1,\ldots, u_n)\right) = \sum_{\substack{I\subseteq[n]\\|I|=k}} (x_{\scriptscriptstyle I})_{_{(1,0)}} =  0\in (\CP_n)_{(1,0)}.\]
It is easy to see that these elements generate the kernel of $\psi_n$. We deduce the isomorphism 
\[ (\CP_n)_{(1,0)} \cong  \C[u_1\,\ldots, u_n]/\langle e_1(u_i), \ldots, e_n(u_i)\rangle = \C[u_1\,\ldots, u_n]/\langle \C[u_1\,\ldots, u_n]^{S_n}_+\rangle
\]
with the classical coinvariant algebra $R_n$. 
The subtorus $\{t_0=1\}\subset (\C^*)^2$ fixes the point $(1,0)\in\A^2$, and acts on the fibre algebra $(\CP_n)_{(1,0)}$, equipping it  with a grading agreeing with the standard grading on the coinvariant algebra. The remaining statements in (iii) are then well known. 

The fibre over a point $(0, s_1)$ is isomorphic to the fibre over a point $(s_0, 0)$ using the ``external'' automorphism of the whole setup which interchanges the coordinate $x_{\scriptscriptstyle I}$ with the coordinate $x_{\scriptscriptstyle [n]\setminus I}$, and the deformation coordinates $s_0$ and $s_1$ (the central reflection of the hypercube). 

Finally for general fibres $(\CP_n)_{(s_0,s_1)}$ with $s_0s_1\neq 0$, we can re-scale to $s_0=s_1=1$. Setting $v_i=(x_{\{i\}})_{(1,1)}\in(\CP_n)_{(1,1)}$, much of the argument of the previous paragraphs goes through, except for the last step where the final relation gets deformed by the parameter $s_1=1$. We get
\[ (\CP_n)_{(1,1)}\cong\C[v_1, \ldots, v_n]\Big/\left\langle \sum_{i=1}^n v_i, \sum_{i<j} v_iv_j, \ldots, \prod_{i=1}^n v_i -1\right\rangle. \]
The latter algebra is the coordinate algebra of the free $S_n$-orbit of the point 
\[(1, \eta, \eta^2, \ldots, \eta^{n-1})\in\A^n = \Spec\C[v_1, \ldots, v_n]\]
for a primitive $n$-th root of unity~$\eta$. Hence this (ungraded) algebra is indeed $S_n$-isomorphic to the regular representation of $S_n$.  

Finally the affine map $\pi_n$ is finite with smooth target and fibres of constant length, hence the inclusion $\C[s_0, s_1]\rightarrow \CP_n$ is indeed flat. 
\end{proof}

\end{document}
